\chapter{Introducción}

\section{El problema}

Guardar cartas es fácil; leerlas no. Una caja de zapatos con mil cartas de
Pokémon no responde ninguna de las tres preguntas que su dueño se hace: qué
tengo, cuánto vale y qué me falta. El almacenamiento no es la parte
interesante del problema —una hoja de cálculo almacena— sino convertir un
montón en algo consultable, valorable e intercambiable.

PokéDex Manager hace eso, y expone los mismos datos por dos caminos: una API
REST que consume la aplicación web, y un \textbf{servidor MCP} al que puede
conectarse cualquier asistente. El chat de la propia aplicación es cliente de
ese servidor MCP, no una segunda ruta hacia los datos. Esa decisión se explica
en el capítulo~\ref{cap:backend} y se ve en los flujos del capítulo~\ref{cap:flujos}.

\section{Qué hace el sistema}

\par\addvspace{8pt}\noindent\begin{pkbox}{colspec={X[1.3,l]X[4,l]}}
  Capacidad & Qué significa técnicamente \\
  Escaneo & Una foto se normaliza, se guarda, se transcribe con un modelo de visión y se resuelve contra el catálogo puntuando cada señal. El usuario confirma. \\
  Catálogo & \pkcode{card\_set}, \pkcode{card} y \pkcode{species} espejados de TCGdex y PokeAPI. La carta es una impresión; la especie es el Pokémon que retrata. \\
  Colección & Grupos homogéneos de copias con condición, idioma, gradeo y coste unitario. \\
  Portafolio & Coste contra valor de mercado, posición por posición, con concentración y serie temporal a partir de una lectura de precio diaria. \\
  Intercambio & Emparejamiento entre coleccionistas, ofertas, contraofertas y un tablón abierto donde una propuesta se publica a nadie en particular. \\
  Mensajería & Hilos directos entre dos coleccionistas, con estado de lectura. \\
  Asistente & Agente sobre 13 herramientas MCP, con memoria de hechos duraderos, respondiendo en español y en streaming. \\
  Compartir & Enlace público revocable con una proyección distinta a la del dueño, más exportación CSV y JSON. \\
\end{pkbox}\par\addvspace{8pt}

\section{Glosario}

\par\addvspace{8pt}\noindent\begin{pkbox}{colspec={X[1.2,l]X[4,l]}}
  Término & Definición en este sistema \\
  Carta (\pkcode{card}) & Una \emph{impresión} concreta: identificador del origen, set, número, rareza, variantes e imagen. \\
  Especie (\pkcode{species}) & El Pokémon retratado. Se une a la carta por el \pkcode{dexId} del origen, nunca por coincidencia de nombres. \\
  Sobrante (\emph{spare}) & Copia excedente de una carta: \pkcode{SUM(quantity) > 1} agrupado por usuario y carta. La primera copia es la colección; lo que sobra es inventario. \\
  Posición & Todas las copias de una misma carta de un usuario colapsadas en una fila, con coste, valor y participación en el portafolio. \\
  Oferta (\pkcode{trade\_offer}) & Propuesta \emph{con destinatario}. Aceptarla registra el acuerdo; no mueve cartas. \\
  Publicación (\pkcode{trade\_listing}) & Propuesta \emph{sin destinatario}. Quien pueda cumplirla la toma, y tomarla es el acuerdo. \\
  MCP & \emph{Model Context Protocol}. El protocolo por el que un asistente externo lee la colección. \\
  JWKS & Conjunto de claves públicas que publica Better Auth y con el que la API verifica los JWT sin llamar de vuelta. \\
  Kubb & Generador que convierte el OpenAPI de la API en tipos, clientes y hooks de React Query. \\
\end{pkbox}\par\addvspace{8pt}
